\begin{center}
\textbf{\large{Tóm tắt}	}
\end{center}
Sự phát triển mạnh mẽ của khoa học vật liệu và công nghệ mô phỏng số đã đưa vật liệu Auxetic – nhóm vật liệu có hệ số Poisson âm – trở thành đối tượng nghiên cứu và ứng dụng đầy tiềm năng trong nhiều lĩnh vực như cơ khí, y sinh và kỹ thuật kết cấu. Tuy nhiên, quá trình thiết kế hình học và điều chỉnh các tham số cấu trúc của vật liệu Auxetic thường đòi hỏi kiến thức chuyên môn sâu cũng như các công cụ mô phỏng phức tạp, gây khó khăn cho người dùng trong việc tiếp cận và thử nghiệm nhanh các phương án thiết kế.

Xuất phát từ thực tế đó, ứng dụng Auxilab được phát triển nhằm cung cấp một nền tảng giao diện web trực quan, cho phép người dùng tương tác trực tiếp và điều chỉnh các tham số thiết kế của cấu trúc Auxetic mà không cần thao tác phức tạp trên các công cụ dòng lệnh hay phần mềm chuyên dụng. Thông qua giao diện web, người dùng có thể dễ dàng thay đổi hình học ô cơ sở, mật độ vật liệu, kích thước đặc trưng và các tham số liên quan, từ đó quan sát sự thay đổi của cấu trúc Auxetic một cách trực quan và tức thời.

Auxilab đóng vai trò là cầu nối giữa người dùng và các thuật toán xử lý – mô phỏng phía sau, giúp đơn giản hóa quy trình thiết kế, tăng khả năng thử nghiệm nhanh nhiều phương án khác nhau, đồng thời tạo tiền đề cho việc tích hợp các bước phân tích, tối ưu hóa và mô phỏng cơ học trong tương lai. Qua đó, hệ thống góp phần nâng cao hiệu quả nghiên cứu và ứng dụng vật liệu Auxetic trong thực tiễn.